\section{Related Work}
\label{sec:related}
This section connects AgentWorkShop to two complementary bodies of work:
plant-facing industrial agents, and the production and automation abstractions
that define context, tool access, and operational constraints. Table~\ref{tab:nearest}
compares selected agent systems; the research gap is limited to these
comparators rather than attributed to the wider literature.

\subsection{Industrial Agents and Plant Integration}
Industrial-agent systems span end-to-end automation control, supervisory
services, programmable logic controller (PLC) interfaces, manufacturing
coordination, and operator-facing machine access. Xia et al. present an
LLM-agent framework for end-to-end industrial automation control, using
structured prompts and event-driven context for production planning and
operation \cite{xia2025control}. Biskupovic et al. study agent-based supervision for
service-oriented industrial systems \cite{biskupovic2025supervision};
Agents4PLC and InstructMPC address PLC code generation and human--LLM
context-aware control, respectively
\cite{liu2026agents4plc,wu2025instructmpc}. Manufacturing studies consider
dynamic resource-capability exploration \cite{lim2025llmenabled}, multi-agent
shop-floor coordination \cite{zhao2026shopfloor}, and digital-twin-supported
fault handling \cite{gill2025llmdt}. Natural-language machine access
\cite{hofmann2025nlc} and real-time LLM control of a robot-operated serial
line \cite{waseem2026pretrained} extend this work to operator interfaces and
robotic-line control. Together, these studies address different operational
contexts, interfaces, and allocations of responsibility between agents,
operators, and controllers.
The comparators were selected for overlap with the framework's integration
scope: end-to-end LLM control, service-oriented industrial agents, and
modular factory skills. Table~I uses common descriptive axes---agent task,
integration interface, and source-reported evaluation. This purposive
comparison is representative rather than exhaustive, and it does not score
system capabilities or compare performance.

\begin{table*}[!t]
\centering
\caption{Selected plant-facing agent systems and their integration settings.
Evidence is described within each source's scope and is not compared as a
performance ranking.}
\label{tab:nearest}
\small
\setlength{\tabcolsep}{4pt}
\begin{tabularx}{\textwidth}{@{}p{2.2cm}XX@{}}
\toprule
System & Integration focus & Source-reported evaluation \\
\midrule
Xia et al. \cite{xia2025control} & End-to-end industrial control with LLM agents, structured prompts, and event-driven context & Proof-of-concept framework; task-specific dataset generation \\
K{\"o}hle et al. \cite{koehle2026llmdriven} & Modular factory skills using OPC UA, Message Queuing Telemetry Transport (MQTT), and component-state tools & Three architectures over nine production challenges \\
Biskupovic et al. \cite{biskupovic2025supervision} & Industrial agents select and reconfigure services from a predefined library using performance and computational-load measures & Online service monitoring and reported adaptability/performance assessment \\
\midrule
AgentWorkShop (this work) & Registered harnesses and shared dispatch; line--product--recipe--run context; semantic I/O and route-specific records & Historical process/engine traces; current two-scene protocol and data integration \\
\bottomrule
\end{tabularx}
\end{table*}

Xia et al. provide the closest overlap in end-to-end industrial-agent scope:
their framework combines an agent system, structured prompting, and
event-driven multilevel context for production planning and operation
\cite{xia2025control}. Biskupovic et al. are a complementary comparison whose
agents select a service from a predefined component library and monitor it for
possible reselection as performance changes
\cite{biskupovic2025supervision}. The integration unit in AgentWorkShop is
different: a channel task is resolved to
production context and a registered harness; the shared dispatcher rechecks
scope, grants, and bindings, and industrial handlers apply route-specific
limits; semantic
nodes route through protocol drivers while command receipts, process
observations, and recipe versions remain separate. This positions the
contribution as a software integration contract around task execution, rather
than a new service-selection or control strategy. The present evidence covers
these components and a two-scene protocol/data path, but not the current
registered-harness-to-process path as one trial. The comparison does not
claim that other systems lack analogous features or establish performance
superiority. Authority and command confirmation are operational services,
consistent with safety-aware decision support \cite{jiang2026agentic}; they
do not replace process protection or establish operator effectiveness.

\subsection{Production Context, Agent Interfaces, and Automation}
Industrial integration also depends on abstractions that are not specific to
agent architectures. MCP standardizes tool discovery and calls with
request-scoped authorization \cite{mcp2024}, whereas
ISA-88 and ISA-95 describe recipes and production operations
\cite{isa88_2010,isa95_2025}. Run-to-run methods connect process observations
to successive production runs \cite{moyne2000r2r}. At the agent level,
AgentSpec and CaMeL address tool behavior and capability-based information
flow \cite{wang2025agentspec,debenedetti2025camel}. Together, these works
provide complementary abstractions for production context, tool interaction,
and information flow. They do not, by themselves, specify the shared
integration architecture presented here.

Industrial security and record requirements further constrain that
architecture: FDA guidance discusses trustworthy records and audit trails
\cite{part11}, IEC~62443-3-3 specifies system-security requirements
\cite{iec62443-3-3}, and DNP3 includes select-before-operate
\cite{ieee1815dnp3}. Control-assurance methods address state evolution through
shielding, control-barrier functions, or Simplex architectures
\cite{alshiekh2018shielding,ames2017cbf,sha2001simplex}. These standards and
methods inform the boundaries within which an industrial-agent framework must
operate; they are not contributions claimed by AgentWorkShop.

Table~\ref{tab:nearest} provides a purposive, descriptive comparison, not a
performance ranking. AgentWorkShop's contribution is implementation-level
integration of registered harnesses, scoped production context, shared
dispatch, semantic process I/O, and route-specific evidence. The evaluation
does not validate per-agent profile adaptation, a new control strategy, or
performance superiority.
